\medskip

L'espace est rapporté à un repère orthonormé \Oijk.

On considère les points suivants:

\begin{center}A(1~;~3~;~0), \quad B$(-1~;~4~;~5)$, \quad C(0~;~1~;~0) \quad et \quad D$(- 2~;~2~;~1)$.\end{center}

\begin{enumerate}
\item Montrer que les points A, B et C déterminent un plan.
\item Montrer que le triangle ABC est rectangle en A.
\item Soit $\Delta$ la droite passant par le point D et de vecteur directeur $\vect{u}\begin{pmatrix}2\\-1\\1\end{pmatrix}$.
	\begin{enumerate}
		\item Démontrer que la droite $\Delta$ est orthogonale au plan (ABC).
		\item Justifier que le plan (ABC) admet pour équation cartésienne:

\[2x - y + z + 1 = 0.\]

		\item Déterminer une représentation paramétrique de la droite $\Delta$.
	\end{enumerate}
\item On appelle H le point de coordonnées $\left(-\frac 23~ ~\frac 43~;~\frac 53\right)$.

Vérifier que H est le projeté orthogonal du point D sur le plan (ABC).
\item On rappelle que le volume d'un tétraèdre est donné par$ V = \dfrac 13 B \times h$, où $B$ est l'aire
 d'une base du tétraèdre et $h$ est sa hauteur relative à cette base.
	\begin{enumerate}
		\item Montrer que DH $=\dfrac{2\sqrt 6}{3}$.
		\item En déduire le volume du tétraèdre ABCD.
	\end{enumerate}
\item On considère la droite $d$ de représentation paramétrique :

\[\left\{\begin{array}{l c r}
x&=&1 - 2k\\
y&=& - 3k\\
z&=&1 + \phantom{3}k\end{array}\right. \: \text{où}\:\: k \:\ \text{décrit}\:\: \R.\]

La droite $d$ et le plan (ABC) sont-ils sécants ou parallèles ?
\end{enumerate}

\bigskip
